\documentclass{article}
\usepackage[margin=1in]{geometry}
\usepackage{amsmath,amssymb}
\usepackage{enumitem}
\usepackage{titlesec}
\usepackage{parskip}
\usepackage{hyperref}
\usepackage{bm}
\usepackage{fancyhdr}
\usepackage{xcolor}
\usepackage{booktabs}
\usepackage{url}
\usepackage{mathtools}

\definecolor{deepblue}{RGB}{0,50,120}
\definecolor{accentblue}{RGB}{30,100,200}
\definecolor{darkgray}{RGB}{80,80,80}

\hypersetup{
    colorlinks=true,
    linkcolor=accentblue,
    urlcolor=accentblue
}

\pagestyle{fancy}
\fancyhf{}
\rhead{\textcolor{deepblue}{\small EEL 6878 --- Spring 2026}}
\lhead{\textcolor{deepblue}{\small Instructor: Tong Wu}}
\rfoot{\thepage}

\titleformat{\section}[block]{\large\bfseries\color{deepblue}}{}{0em}{}[\titlerule]
\titleformat{\subsection}[block]{\normalsize\bfseries\color{accentblue}}{}{0em}{}

% Custom question counter
\newcounter{questioncounter}
\newcommand{\question}[1]{%
  \stepcounter{questioncounter}%
  \vspace{1em}%
  \noindent\textbf{\textcolor{deepblue}{\large Question \thequestioncounter.} #1}%
  \vspace{0.3em}%
}

\DeclarePairedDelimiter\abs{\lvert}{\rvert}

\begin{document}

%% ============================================================
%%  TITLE PAGE / HEADER
%% ============================================================
\begin{center}
  {\Large \textbf{EEL 6878: Modeling and Artificial Intelligence}}\\[0.3em]
  {\large Spring 2026}\\[0.2em]
  {\small \textcolor{darkgray}{Instructor: Tong Wu}}\\[1em]
  \rule{\linewidth}{0.8pt}\\[0.4em]
  {\LARGE \textbf{\textcolor{deepblue}{Homework: Graph Signal Processing \& Graph Neural Networks}}}\\[0.3em]
  \rule{\linewidth}{0.8pt}
\end{center}

\vspace{1em}

\tableofcontents

\newpage

%% ============================================================
%%  PART 1: GRAPH SIGNAL PROCESSING
%% ============================================================
\section{Graph Signal Processing}

In this section, you will use the existing package \texttt{pygsp}\footnote{\url{https://pygsp.readthedocs.io/en/stable/}} to complete tasks related to graph signal processing.

\setcounter{questioncounter}{0}

\question{\textbf{Spectral Analysis of Graph Signals}}

The spectral content of a signal indicates whether it is low-pass, band-pass, or high-pass. Intuition transfers from classical Fourier analysis to the graph setting.

\begin{enumerate}[leftmargin=1.5em]
    \item Build a sensor graph using:
    \begin{center}
        \texttt{G = pygsp.graphs.Sensor(N=100, seed=1)}
    \end{center}

    \item Compute its Graph Fourier basis.

    \item Create two kinds of graph signals by following these steps:
    \begin{enumerate}[leftmargin=1.5em]
        \item Generate random signals of size $1$ for each node in the graph according to the standard normal distribution.
        \item Generate two different filter banks of Heat kernels\footnote{\url{https://pygsp.readthedocs.io/en/latest/reference/filters.html}} characterized by their respective parameters: $\mathtt{scale} = 0$ for the first filter and $\mathtt{scale} = 10$ for the second filter. These parameters alter the filter's behavior, leading to different spectral characteristics of the processed signals.
        \item Apply the defined filters to the randomly generated signals. This process will create two distinct types of graph signals: one demonstrating low-pass characteristics and the other high-pass characteristics.
    \end{enumerate}

    \item Calculate the Graph Fourier Transform (GFT) for both the low-pass and high-pass graph signals to obtain the signals in the Graph Fourier domain $\hat{\bm{x}}$.

    \item Plot the graph spectral content of $\hat{\bm{x}}$, i.e., $\abs{\hat{\bm{x}}}$, for both signals.

    \item Evaluate the smoothness of the two graph signals, defined as:
    \[
        S(\bm{x}) = \bm{x}^\top \mathbf{L}\, \bm{x}
    \]
    where $\mathbf{L}$ is the graph Laplacian. Compare and discuss the smoothness values of the two signals.

    \item Is the graph signal generated by $\mathtt{scale}=10$ a high-pass or a low-pass graph signal? Justify your answer.
\end{enumerate}

\vspace{1em}

\question{\textbf{Synthetic and Real-World Networks}}

\begin{enumerate}[leftmargin=1.5em]
    \item Generate (or load) the following \textit{undirected} graphs and plot them:
    \begin{enumerate}[leftmargin=1.5em]
        \item An Erd\H{o}s-R\'{e}nyi graph with $N = 50$ and $p = 0.1$ using \texttt{pygsp}.
        \item A Ring graph with $N = 50$ using \texttt{pygsp}.
        \item The Cora\footnote{The Cora dataset consists of 2{,}708 scientific publications classified into one of seven classes and 5{,}429 citation links. Each publication is described by a $0/1$-valued word vector of length 1{,}433 indicating word presence.} citation network from PyTorch Geometric.\footnote{\url{https://pytorch-geometric.readthedocs.io/en/latest/modules/datasets.html}}
    \end{enumerate}

    \item Since we are considering undirected graphs, the Laplacian matrix is real symmetric and thus admits an eigen-decomposition:
    \[
        \mathbf{L} = \mathbf{U}\mathbf{\Lambda}\mathbf{U}^\top,
    \]
    where $\mathbf{U} = (\mathbf{u}_1, \ldots, \mathbf{u}_N) \in \mathbb{R}^{N \times N}$ is the matrix of orthonormal eigenvectors and $\mathbf{\Lambda} = \mathrm{diag}(\lambda_1, \ldots, \lambda_N)$ with sorted eigenvalues:
    \[
        \lambda_1 \leq \lambda_2 \leq \cdots \leq \lambda_N.
    \]
    \begin{itemize}[leftmargin=1.5em]
        \item Plot the graph frequencies (eigenvalues) of the Erd\H{o}s-R\'{e}nyi, Ring, and Cora graphs, from smallest to largest.
        \item Plot the eigenvectors $\mathbf{u}_2$, $\mathbf{u}_3$, and $\mathbf{u}_6$ overlaid on the Cora graph using a node colormap.\footnote{See \url{https://networkx.org/documentation/stable/auto_examples/drawing/plot_node_colormap.html} for guidance.}
    \end{itemize}

    \item Compute the covariance matrix of the Cora node feature data. Use t-SNE to obtain two principal components for visualization, then apply $K$-means clustering with $K = 7$ classes. Report:
    \begin{itemize}[leftmargin=1.5em]
        \item The Adjusted Rand Index (ARI) of the resulting labels.
        \item The Silhouette score of the resulting labels.
    \end{itemize}

    \item Let's define a low-pass filter in the spectral domain:
    \[
        \tilde{h}(\lambda) = \frac{1}{1 + \alpha\lambda}.
    \]
    Given a noisy version $\bm{x}_{\mathrm{noisy}}$ of a smooth signal, one can denoise it as:
    \[
        \bm{x}_{\mathrm{denoised}} = \mathbf{U}\,\tilde{h}(\mathbf{\Lambda})\,\mathbf{U}^\top\, \bm{x}_{\mathrm{noisy}}.
    \]
    \begin{enumerate}[leftmargin=1.5em]
        \item Consider the Ring graph from Part (a) and define the following graph signal:
        \[
            \bm{x} = \sin(i) + \cos(2i) + 0.5\cos(30i),
        \]
        where $i = \texttt{np.linspace}(-\pi, \pi, N)$. Add white noise drawn from a Uniform distribution on $[-0.4,\, 0.4]$ to the signal.
        \item Plot the noisy graph signal.
        \item Compute the signal in the Graph Fourier domain $\hat{\bm{x}}$ and plot its spectral content $\abs{\hat{\bm{x}}}$.
        \item Define the low-pass filter with $\alpha = 2$ using:
        \begin{center}
            \texttt{g = pygsp.filters.Filter(G, h)}
        \end{center}
        \item Apply $g$ to denoise the signal:
        \begin{center}
            \texttt{x\_denoised = g.filter(x, method='exact')}
        \end{center}
        \item Plot the denoised graph signal and its graph spectral content. Compare with the noisy version and discuss the effect of the low-pass filter.
    \end{enumerate}
\end{enumerate}

\newpage

%% ============================================================
%%  PART 2: GRAPH NEURAL NETWORKS
%% ============================================================
\section{Graph Neural Networks}

\setcounter{questioncounter}{2}

\question{\textbf{Node Classification with Graph Neural Networks}}

In this problem, you will use PyTorch Geometric to build graph neural networks for node classification on the Cora dataset.\footnote{A tutorial on GNN implementation is available at \url{https://pytorch-geometric.readthedocs.io/en/latest/notes/introduction.html}.}

\begin{enumerate}[leftmargin=1.5em]
    \item \textbf{Model Architecture.} Build a two-layer Chebyshev Graph Neural Network (ChebNet) using PyTorch Geometric's \texttt{ChebConv} operator with the following specifications:
    \begin{itemize}[leftmargin=1.5em]
        \item \textbf{Layer 1:} Input channels = number of node features in the dataset; output channels = 40; Chebyshev order $K = 5$.
        \item \textbf{Layer 2:} Input channels = 40; output channels = number of label classes in the dataset.
        \item \textbf{Activation functions:} Use \texttt{relu} after the first layer and \texttt{log\_softmax} at the output layer to produce predicted class probabilities.
    \end{itemize}

    \item \textbf{Training.} Train your GNN using the following setup:
    \begin{itemize}[leftmargin=1.5em]
        \item Optimizer: \texttt{Adam} with learning rate $0.01$ and weight decay $5 \times 10^{-4}$.\footnote{You may adjust the learning rate and weight decay as needed for your model.}
        \item Loss function: \texttt{torch.nn.functional.nll\_loss} (negative log-likelihood loss), or any other suitable loss for your model.
        \item Plot the training loss as a function of the epoch number.
    \end{itemize}

    \item \textbf{Evaluation and Visualization.} Test your trained GNN and do the following:
    \begin{itemize}[leftmargin=1.5em]
        \item Report the classification accuracy on the test set.
        \item Visualize the Cora graph with node colors indicating (i) ground truth labels and (ii) predicted labels. Discuss any notable differences.
    \end{itemize}
\end{enumerate}

%% ============================================================
%%  BIBLIOGRAPHY
%% ============================================================
\newpage
\begin{thebibliography}{9}

\bibitem{pygsp}
N.~Perraudin et al.,
``GSPBOX: A toolbox for signal processing on graphs,''
\textit{arXiv preprint arXiv:1408.5781}, 2014.
\url{https://pygsp.readthedocs.io/en/stable/}

\bibitem{cora}
A.~K.~McCallum, K.~Nigam, J.~Rennie, and K.~Seymore,
``Automating the Construction of Internet Portals with Machine Learning,''
\textit{Information Retrieval}, vol.~3, pp.~127--163, 2000.

\bibitem{pyg}
M.~Fey and J.~E.~Lenssen,
``Fast Graph Representation Learning with PyTorch Geometric,''
in \textit{ICLR Workshop on Representation Learning on Graphs and Manifolds}, 2019.

\end{thebibliography}

\end{document}
